\chapter{Direct limits of topological spaces and \texorpdfstring{$k_\omega$}{k-omega}-spaces}
\label{chap:topology}

Direct limits of ascending sequences of topological spaces, groups and vector spaces, after
\cite[Appendix~A.7]{GN26}. They are the topological side of the locally convex direct limits of
Chapter~\ref{chap:directsums}. The main points are that the direct limit of a compactifying sequence of
Hausdorff groups is a topological group, and that $k_\omega$-spaces (direct limits of ascending
sequences of compact spaces) are stable under direct limits and finite products. Nothing in
this chapter depends on the calculus.

\section{Direct sequences and direct limits}

\begin{definition}[Direct sequences {\cite[Definitions~A.7.1, A.7.5, Remark~A.7.15]{GN26}}]
  \label{def:top-dl}
  \lean{LocallyConvex.Analysis.TopDirectLimit, LocallyConvex.Analysis.IsDirectSeq,
    LocallyConvex.Analysis.IsDirectSeq.IsStrict, LocallyConvex.Analysis.IsDirectSeq.IsCompactifying,
    LocallyConvex.Analysis.IsDirectSeq.CompactlyRegular}
  \leanok
  Let $X_0 \subseteq X_1 \subseteq \cdots$ be topological spaces covering a set $X$, with continuous
  inclusions. The \emph{direct limit topology} on $X$ is the final topology of the inclusions. The
  sequence is \emph{strict} if each $X_n$ carries the topology induced by $X_{n+1}$, and
  \emph{compactifying} if every point of $X_n$ has a neighbourhood whose closure in some $X_m$ is
  compact. It is \emph{compactly regular} if every compact subset of the limit lies, and is
  compact, in some $X_n$.
\end{definition}

\begin{theorem}[Basic properties {\cite[Remark~A.7.2, Lemmas~A.7.3, A.7.6, Remark~A.7.15]{GN26}}]
  \label{thm:top-dl-basic}
  \lean{LocallyConvex.Analysis.TopDirectLimit.isOpen_iff,
    LocallyConvex.Analysis.TopDirectLimit.continuous_iff,
    LocallyConvex.Analysis.TopDirectLimit.eq_topology_iff,
    LocallyConvex.Analysis.TopDirectLimit.topology_eq_comp,
    LocallyConvex.Analysis.TopDirectLimit.continuous_map,
    LocallyConvex.Analysis.IsDirectSeq.isOpen_iUnion,
    LocallyConvex.Analysis.IsDirectSeq.exists_subset_of_isCompact,
    LocallyConvex.Analysis.IsDirectSeq.IsStrict.isEmbedding_incl,
    LocallyConvex.Analysis.IsDirectSeq.IsStrict.compactlyRegular,
    LocallyConvex.Analysis.IsDirectSeq.isCompactifying_of_locallyCompactSpace,
    LocallyConvex.Analysis.IsDirectSeq.isStrict_of_compactSpace,
    LocallyConvex.Analysis.TopDirectLimit.topology_subtype_eq,
    LocallyConvex.Analysis.IsDirectSeq.topology_iUnion_eq}
  \leanok
  \uses{def:top-dl}
  A set is open (closed) in the direct limit iff its traces on all steps are; a map on the limit is
  continuous iff its restrictions to the steps are, and this characterizes the topology; cofinal
  subsequences give the same limit; ascending unions of open subsets of the steps are open. If the
  steps are $T_1$, so is the limit, and compact sets lie in some step. Strict sequences have
  embedded steps and, for $T_1$ steps, are compactly regular. Locally compact Hausdorff steps are
  compactifying; compact Hausdorff steps are strict. Open and closed subsets carry the direct limit
  topology of their traces, and so do ascending unions $\bigcup_n V_n$ of open subsets
  $V_n \subseteq X_n$, with respect to the $V_n$.
\end{theorem}

\begin{proof}
  \leanok
  As in \cite{GN26}. For compact sets, points $x_n \in K \setminus X_n$ would form an infinite
  closed discrete subset of $K$. For strict sequences, open subsets of $X_m$ are extended
  recursively to open subsets of the later steps.
\end{proof}

\begin{theorem}[Compactifying sequences {\cite[Lemmas~A.7.4, A.7.6(h), A.7.9,
  Proposition~A.7.20]{GN26}}]
  \label{thm:top-dl-compactifying}
  \lean{LocallyConvex.Analysis.IsDirectSeq.IsCompactifying.exists_isOpen,
    LocallyConvex.Analysis.IsDirectSeq.IsCompactifying.exists_disjoint,
    LocallyConvex.Analysis.IsDirectSeq.IsCompactifying.t2Space,
    LocallyConvex.Analysis.IsDirectSeq.IsCompactifying.isOpen_prod_iff,
    LocallyConvex.Analysis.IsDirectSeq.IsCompactifying.continuous_prod_iff,
    LocallyConvex.Analysis.IsDirectSeq.IsCompactifying.compactlyRegular,
    LocallyConvex.Analysis.IsDirectSeq.IsCompactifying.isLocallyKOmegaSpace,
    LocallyConvex.Analysis.isKOmegaSpace_iff_exists_topology_eq,
    LocallyConvex.Analysis.IsKOmegaSpace.isLocallyKOmegaSpace,
    LocallyConvex.Analysis.isKOmegaSpace_of_sigmaCompactSpace,
    LocallyConvex.Analysis.isLocallyKOmegaSpace_of_locallyCompactSpace,
    LocallyConvex.Analysis.IsLocallyKOmegaSpace.compactlyGeneratedSpace}
  \leanok
  \uses{thm:top-dl-basic}
  Let $X_0 \subseteq X_1 \subseteq \cdots$ and $Y_0 \subseteq Y_1 \subseteq \cdots$ be
  compactifying sequences of Hausdorff spaces. Then $\varinjlim X_n$ is Hausdorff, compactly
  regular and locally $k_\omega$, and $\varinjlim X_n \times \varinjlim Y_n =
  \varinjlim (X_n \times Y_n)$ as topological spaces. A Hausdorff space is $k_\omega$ iff it is
  the direct limit of an ascending sequence of compact subspaces; $k_\omega$-spaces and locally
  compact spaces are locally $k_\omega$, and locally $k_\omega$-spaces are $k$-spaces
  \cite[Remark~A.7.19]{GN26}.
\end{theorem}

\begin{proof}
  \leanok
  Compact subsets of a step have open neighbourhoods with compact closure in a later step, inside
  any prescribed open set. Iterating gives chains $K_j \subseteq V_j \subseteq K_{j+1}$, with $V_j$
  open and $K_j$ compact in increasing steps, whose unions are open: two disjoint chains separate
  points, and products of chains fit inside open subsets of the product. A set meeting every
  $K_j$ in a closed set is closed in $\bigcup_j V_j$, so compact subsets of $\bigcup_j V_j$ lie in
  some $K_j$ (otherwise they contain an infinite closed discrete set); covering a compact set by
  finitely many compact pieces of such unions gives compact regularity.
\end{proof}

\begin{theorem}[Groups and vector spaces {\cite[Proposition~A.7.11, Corollary~A.7.12,
  Example~A.7.13]{GN26}}]
  \label{thm:top-dl-groups}
  \lean{LocallyConvex.Analysis.TopDirectLimit.continuous_inv',
    LocallyConvex.Analysis.TopDirectLimit.continuous_mul_left',
    LocallyConvex.Analysis.TopDirectLimit.continuous_mul_right',
    LocallyConvex.Analysis.IsDirectSeq.IsCompactifying.isTopologicalGroup,
    LocallyConvex.Analysis.IsDirectSeq.IsCompactifying.nhds_one_hasBasis,
    LocallyConvex.Analysis.IsDirectSeq.isTopologicalGroup_of_locallyCompactSpace,
    LocallyConvex.Analysis.IsDirectSeq.IsCompactifying.continuousSMul,
    LocallyConvex.Analysis.IsDirectSeq.IsCompactifying.locallyConvexSpace,
    LocallyConvex.Analysis.IsDirectSeq.IsCompactifying.topology_eq_lcDirectLimit}
  \leanok
  \uses{thm:top-dl-compactifying, thm:ds-limit}
  For an ascending sequence of topological groups with continuous inclusions, inversion and
  translations are continuous in the direct limit topology. If the sequence is compactifying and
  the steps are Hausdorff, the limit is a topological group, with the sets
  $\bigcup_n U_0 U_1 \cdots U_n$ ($U_n$ open identity neighbourhoods of the steps) as a basis of
  identity neighbourhoods. Compactifying direct limits of Hausdorff (locally convex) topological
  vector spaces are (locally convex) topological vector spaces, and the direct limit topology is
  then the locally convex direct limit topology.
\end{theorem}

\begin{proof}
  \leanok
  Multiplication and scalar multiplication are continuous by the product theorem, the field being
  the constant compactifying sequence. In the basis the $U_n$ can be taken convex, and ascending
  unions of convex sets are convex. A locally convex topology making the inclusions continuous is
  coarser than the locally convex direct limit topology, by its universal property.
\end{proof}

\begin{theorem}[Further properties {\cite[Remark~A.7.7, Example~A.7.17, Remark~A.7.21,
  Lemma~A.7.22]{GN26}, \cite[Section~4]{GGH10}}]
  \label{thm:top-dl-komega}
  \lean{LocallyConvex.Analysis.IsDirectSeq.exists_nested_nhds,
    LocallyConvex.Analysis.IsDirectSeq.IsCompactifying.isCompactifying_unionSet,
    LocallyConvex.Analysis.compactlyGeneratedSpace_iff_topology_eq,
    LocallyConvex.Analysis.IsDirectSeq.t2Space_isKOmegaSpace,
    LocallyConvex.Analysis.IsDirectSeq.isLocallyKOmegaSpace,
    LocallyConvex.Analysis.IsKOmegaSpace.prod,
    LocallyConvex.Analysis.IsLocallyKOmegaSpace.prod,
    LocallyConvex.Analysis.IsDirectSeq.isOpen_prod_iff,
    LocallyConvex.Analysis.IsDirectSeq.isOpen_prod_iff_of_isKOmegaSpace}
  \leanok
  \uses{thm:top-dl-compactifying}
  For locally compact Hausdorff steps, a point of an open set has nested open neighbourhoods
  $V_n \subseteq X_n$ with compact closures inside the set and $\overline{V_n} \subseteq V_{n+1}$.
  The sets $V_n$ of Lemma~A.7.6(b) form a compactifying sequence if the $X_n$ do. A Hausdorff
  space is a $k$-space iff it is the direct limit of its compact subsets. Direct limits of
  (locally) $k_\omega$-spaces are (locally) $k_\omega$ and compactly regular, finite products of
  (locally) $k_\omega$-spaces are (locally) $k_\omega$, and
  $\varinjlim X_n \times \varinjlim Y_n = \varinjlim (X_n \times Y_n)$ for locally $k_\omega$
  steps.
\end{theorem}

\begin{proof}
  \leanok
  As in \cite{GN26} and \cite{GGH10}. For $k_\omega$ steps, the diagonal of cumulative
  $k_\omega$-sequences is an ascending sequence of compact spaces with the same direct limit. For
  locally $k_\omega$ steps, open $k_\omega$ sets $W_n \subseteq X_n$ with $W_n \subseteq W_{n+1}$
  (countable unions of open $k_\omega$ sets are $k_\omega$) have union $\varinjlim W_n$. The
  product of two such limits is a $k$-space, and its compact subsets lie in products of compact
  subsets of the steps.
\end{proof}

\begin{theorem}[The finite topology {\cite[Examples~A.7.14, A.7.16]{GN26}}]
  \label{thm:top-dl-finite}
  \lean{LocallyConvex.Analysis.FiniteTop, LocallyConvex.Analysis.FiniteTop.continuous_linearMap,
    LocallyConvex.Analysis.FiniteTop.topology_le,
    LocallyConvex.Analysis.FiniteTop.topology_eq_seq,
    LocallyConvex.Analysis.FiniteTop.locallyConvexSpace,
    LocallyConvex.Analysis.FiniteTop.continuousLinearEquivFinsupp}
  \leanok
  \uses{thm:top-dl-groups}
  The direct limit topology of the finite-dimensional subspaces of a real vector space $E$ is finer
  than every vector topology, and linear maps on it are continuous. If $E$ has countable
  dimension, it is a Hausdorff locally convex vector topology, hence the finest vector topology
  and the finest locally convex topology, and $E \cong \mathbb{R}^{(\mathbb{N})}$.
\end{theorem}

\begin{proof}
  \leanok
  Linear maps on finite-dimensional spaces are continuous. For countable dimension, the spans of
  the first $n$ basis vectors are cofinal and locally compact, so the limit is a locally convex
  vector topology by Corollary~A.7.12.
\end{proof}
